\section*{Chapitre \ref{ch:rc:convexity-adjustments-and-constant-maturity-products} --- Ajustements de convexité et produits à maturité constante}

\begin{solution}{exo:rc:convexity-adjustments-and-constant-maturity-products:1}
$\delta\sigma_N^2T/(1+\delta F) = 0.25\times0.008^2\times4/(1+0.25\times0.025) =
0.64$ points de base.
\end{solution}

\begin{solution}{exo:rc:convexity-adjustments-and-constant-maturity-products:2}
Versé peu après l'observation, le numéraire de la mesure de paiement $P(T,T_p)$ est
presque égal à un tandis que l'annuité baisse quand les taux montent, de sorte que $P/A$ croît avec $S$~:
$a_1>0$, et la mesure de paiement pondère davantage les états de taux élevés que ne le fait
la probabilité annuité. Avec des taux positifs, $a_0 + a_1S(0) = P(0,T_p)/A(0) >
1/\sum\delta_i = a_0$, d'où $a_1 > 0$.
\end{solution}

\begin{solution}{exo:rc:convexity-adjustments-and-constant-maturity-products:3}
$a_0 = 1/(10\times1.0139) = 0.09863$.
\end{solution}

\begin{solution}{exo:rc:convexity-adjustments-and-constant-maturity-products:4}
Avec un smile plat, l'ajustement est proportionnel à $\sigma_N^2$~:
$11.73\times1.2^2 = 16.89$ points de base.
\end{solution}

\begin{solution}{exo:rc:convexity-adjustments-and-constant-maturity-products:5}
Environ 29~\%~: les swaptions situées à plus de deux points de pourcentage hors de la monnaie portent
près d'un tiers de l'ajustement à dix ans. Les ailes du cube, que le marché des swaptions
cote peu, doivent être valorisées avec soin, bornées par des densités
(chapitre~5), et sont souvent calibrées sur les spreads CMS négociés.
\end{solution}

\begin{solution}{exo:rc:convexity-adjustments-and-constant-maturity-products:6}
$+9$ points de base~: quand les taux en euros montent, l'euro tend à baisser, de sorte que chaque
point de base en euros vaut moins de dollars dans les états où les taux sont élevés~; la mesure
de paiement en dollars décale vers le haut l'espérance du taux.
\end{solution}

\begin{solution}{exo:rc:convexity-adjustments-and-constant-maturity-products:7}
$-19.24$ points de base. Versé à la fin du swap, $P(T,T_p)/A(T)$ baisse quand les taux
montent (le facteur d'actualisation long baisse plus vite que l'annuité), donc $a_1<0$ et
la mesure de paiement pondère davantage les états de taux bas~: l'ajustement est négatif.
\end{solution}

\begin{solution}{exo:rc:convexity-adjustments-and-constant-maturity-products:8}
Chaque coupon CMS est un taux de swap versé une seule fois, à une date qui ne correspond pas au
numéraire de l'annuité sous lequel le taux de swap forward est une martingale~; sa valeur
est le forward plus un ajustement de convexité qui vaut un strip de swaptions (32
points de base pour le taux à dix ans dans dix ans ici). Valorisés aux forwards, les
coupons sont sous-évalués, et une couverture par des swaps forward laisse le livre vendeur
de volatilité~; il faut couvrir l'ajustement par des swaptions sur l'ensemble des strikes.
\end{solution}

\begin{solution}{pb:rc:convexity-adjustments-and-constant-maturity-products:1}
\textbf{1.} Une \omterm{def:rc:convexity-adjustments-and-constant-maturity-products:spread}{option sur spread CMS} de strike zéro sur le taux à dix ans moins le taux à deux ans,
versée un an après l'observation, multipliée par $p$.
\textbf{2.} Chacun est un taux de swap versé une seule fois, un an après l'observation~: chacun
requiert un ajustement CMS~; celui du taux à dix ans est plus grand, son annuité étant plus
sensible aux taux et sa volatilité plus élevée par unité de $a_1$.
\textbf{3.} 0.17118 par unité de notionnel.
\textbf{4.} 0.05785.
\textbf{5.} $0.17118/0.05785 = 2.96$.
\textbf{6.} 2.44, 3.10 et 4.47.
\textbf{7.} Corrélation plus élevée, volatilité du spread plus faible, floors sur spread moins chers
de strike zéro (le forward du spread étant positif)~: chaque unité de coupon
coûte moins, de sorte que davantage d'unités entrent dans le même budget.
\textbf{8.} Il est vendeur des options sur spread qu'il verse, donc vendeur de volatilité du spread et
acheteur de corrélation~: il perd si la corrélation baisse.
\textbf{9.} 45.1 points de base.
\textbf{10.} Les taux de swap et les rendements du Trésor diffèrent par les spreads de swap, qui évoluent
selon leur propre dynamique~; la corrélation qui compte est celle sous la probabilité de valorisation
et sur l'horizon de la note, et non les variations historiques quotidiennes.
\textbf{11.} Des swaptions payeuses et receveuses hors de la monnaie sur le swap sous-jacent de chaque
taux, sur l'ensemble des strikes, dans les poids de réplication.
\textbf{12.} Les ajustements CMS augmentent, ce qui relève les forwards CMS des deux taux~; l'effet
sur le spread dépend de celui qui monte le plus, et les options sur spread profitent
de la volatilité plus élevée de chaque jambe.
\textbf{13.} Les deux~: le niveau par les ajustements CMS et l'application de l'annuité,
la pente directement par le forward du spread.
\textbf{14.} Les coupons tombent à zéro~: le floor à zéro protège l'investisseur contre
des coupons négatifs, mais non contre le fait de ne rien recevoir pendant des années.
\textbf{15.} De nombreux investisseurs achètent la même transaction auprès de quelques dealers, qui finissent vendeurs
de la même corrélation et de la même volatilité du spread et ne peuvent la couvrir
qu'entre eux~; le livre grossit avec les émissions et se réduit difficilement.
\textbf{16.} Inférieur à 2.96~: la banque facture ses coûts de couverture, sa réserve de corrélation
et son financement~; une marge de quelques dixièmes de la participation est usuelle pour
une corrélation peu liquide.
\textbf{17.} À partir d'estimations historiques sur plusieurs fenêtres et régimes, à partir des
corrélations implicites de toute option sur spread négociée, avec une réserve pour l'écart
entre elles~; il faut la soumettre à un stress dans le rapport de risque.
\textbf{18.} Que la note ne verse que si la courbe est plus pentue que zéro à chaque
fixing~; une courbe plate ou inversée pendant des années ne verse rien, et les forwards
intègrent déjà le pentement attendu.
\textbf{19.} Résultat~: \emph{la note steepener} peut verser 2.96 fois le spread dix ans
moins deux ans à une corrélation de 0.769, entre 2.44 et 4.47 à mesure que la
corrélation passe de 0.6 à 0.95.
\textbf{20.} Un strip d'options sur la pente de la courbe, financé en renonçant
à un coupon fixe.
\end{solution}

\begin{solution}{iq:rc:convexity-adjustments-and-constant-maturity-products:1}
La différence entre l'espérance d'un taux sous la probabilité de son paiement effectif
et son forward, provoquée par la corrélation entre le taux et le rapport
des numéraires. Exemples~: contrats à terme contre forwards (appels de marge quotidiens),
un taux versé à sa date de fixing (calendrier), un taux CMS, un taux versé dans une autre
devise (quanto).
\iqlookfor{l'explication par le changement de mesure et plusieurs exemples.}
\end{solution}

\begin{solution}{iq:rc:convexity-adjustments-and-constant-maturity-products:2}
Passer de la probabilité forward $T_k$ à la probabilité forward $T_{k-1}$, de densité
$(1+\delta F_k(T_{k-1}))/(1+\delta F_k(0))$~; alors $\E^{T_{k-1}}[F] =
F_0+\delta\Var^{T_k}(F)/(1+\delta F_0)$, soit $\delta\sigma_N^2T/(1+\delta F_0)$ dans le
modèle normal.
\iqlookfor{la densité de Radon--Nikodym entre probabilités forward.}
\end{solution}

\begin{solution}{iq:rc:convexity-adjustments-and-constant-maturity-products:3}
Écrire le payoff sous la probabilité annuité~: valeur $=A(0)\E^A[(S-K)^+P(T,T_p)/A(T)]$~;
transformer $P/A$ en $a_0+a_1S$ (\omterm{def:rc:convexity-adjustments-and-constant-maturity-products:tsr}{modèle linéaire du taux swap terminal})~; développer
$(S-K)^+(a_0+a_1S) = (a_0+a_1K)(S-K)^+ + 2a_1\int_K^\infty(S-L)^+dL$ et évaluer
chaque swaption payeuse sur le smile.
\iqlookfor{changement vers l'annuité et réplication statique.}
\end{solution}

\begin{solution}{iq:rc:convexity-adjustments-and-constant-maturity-products:4}
Vers le bas~: les coupons CMS que l'on doit valent le forward plus une bande de
swaptions hors de la monnaie, si bien qu'on est vendeur de volatilité, ailes comprises~;
les strikes proches et éloignés du forward comptent, près d'un tiers de
l'ajustement provenant de strikes situés à plus de deux points de pourcentage.
\iqlookfor{position vendeuse en vega via la réplication, et les ailes.}
\end{solution}

\begin{solution}{iq:rc:convexity-adjustments-and-constant-maturity-products:5}
Valoriser à partir d'une plage d'estimations historiques et des corrélations implicites
issues d'options sur spread négociées ou de consultations de dealers, avec une réserve pour l'incertitude~;
couvrir l'exposition à la corrélation par des produits de spread de sens contraire lorsque
c'est possible, sinon la gérer dans des limites et la soumettre à des tests de résistance~; présenter
séparément la sensibilité à la corrélation.
\iqlookfor{valorisation avec réserves, et franchise sur le risque impossible à couvrir.}
\end{solution}

\begin{solution}{iq:rc:convexity-adjustments-and-constant-maturity-products:6}
Largeur de la grille (suffisamment d'écarts types, et aucun strike sous le plancher du modèle),
pas resserré près du forward là où l'intégrande est pentue, selon que le strike et le
forward sont des nœuds de la grille, le comportement du smile sur les ailes (extrapolation,
NaN), et la convergence lorsque la grille est raffinée~; comparer à la forme fermée
sous un smile plat.
\iqlookfor{rigueur de l'intégration numérique et une référence.}
\end{solution}
